% Ad Atticum 7.20 (ad-atticum-07-20) --- English translation
% Translated by Claude (Anthropic) from the Latin (Perseus).
% Style guide: STYLE.md.

\ciceroLetterOpener{CICERO ATTICO salutem}

\ciceroSection{1} The moment itself is making me terse. I have given up on peace; our side is conducting no war. For do not suppose that anything is worth less than these consuls of ours --- in hope of hearing something from whom, and of getting some sense of our preparations, I came to Capua on the day before the Nones in a downpour, as I had been ordered. But they had not yet arrived, and were going to arrive empty-handed and unprepared. Gnaeus, meanwhile, was said to be at Luceria and to be visiting the cohorts of the Appian legions --- not the steadiest of troops. But they report that the other is hurtling on and is at the very gates, not to come to grips (with whom, after all?) but to cut off any retreat.

\ciceroSection{2} For my part, if I stay in Italy, \textit{[Greek: I will die with them too]} --- nor am I consulting you about that; but if outside, what do I do? To stay, I am urged by winter, by my lictors, by leaders who are improvident and careless; to flee, by friendship with Gnaeus, by the cause of the loyalists, by the disgrace of joining hands with a tyrant --- and which one he is going to imitate, Phalaris or Pisistratus, is still uncertain. I should be glad if you would untangle these things and help me with your counsel; though I suppose you yourselves at Rome are now in the thick of it, still --- so far as you can. If I learn anything new here today, you will know it; for the consuls will soon be here for their own Nones. I shall be looking out for your letters every day; and to this one you will write back when you can. The women and the young Ciceros I have left at the Formian villa.
